\begin{table}[t]
\caption{Total Effect Without Topic Adjustment (Probabilistic Time)} 
\fontsize{12.0pt}{14.0pt}\selectfont
\begin{tabular*}{\linewidth}{@{\extracolsep{\fill}}lcccccc}
\toprule
  & H1 Headline‡ & H1 Total‡ & H2 Headline† & H2 Total† & H3a Headline† & H3a Total† \\ 
\midrule\addlinespace[2.5pt]
Invasion & -0.217*** & -0.240*** & -0.253*** & -0.296*** & -0.068 & -0.053 \\ 
 & (0.036) & (0.035) & (0.036) & (0.031) & (0.080) & (0.076) \\ 
State-Owned x Treat &  &  & 0.080 & 0.096 &  &  \\ 
 &  &  & (0.063) & (0.058) &  &  \\ 
Russia-Aligned x Treat &  &  &  &  & -0.19 & -0.22+ \\ 
{} & {} & {} & {} & {} & {(0.11)} & {(0.12)} \\ 
R2 &  &  & 0.045 & 0.022 & 0.025 & 0.010 \\ 
R2 Adj. &  &  & 0.045 & 0.021 & 0.024 & 0.009 \\ 
Topic FEs & X & - & X & - & X & - \\ 
Source Z-Score & X & X & X & X & X & X \\ 
N (within bandwidth) & 45,366 & 45,398 & 40,710 & 45,304 & 7,164 & 6,999 \\ 
Clusters & — & — & 60 & 60 & 13 & 13 \\ 
\bottomrule
\end{tabular*}
\begin{minipage}{\linewidth}
\vspace{.05em}
+ p < 0.1, * p < 0.05, ** p < 0.01, *** p < 0.001\\
Total-effect columns remove the topic adjustment from the corresponding headline
specification and are a distinct estimand, not a robustness check: they include
shifts in the topic composition of coverage rather than holding it fixed.
Sample, source-level Z-Score outcome and clustering are unchanged; all articles
are used, military topics included. The MSE-optimal bandwidth is re-selected
for each total-effect column with topic removed from the residualization, so
the bandwidth remains optimal for the specification actually estimated. H1 is an
rdrobust fit; the H2 and H3a columns are the state-ownership and Russia-aligned
interactions respectively. † Standard errors computed via wild cluster
restricted (WCR) bootstrap (Cameron, Gelbach \& Miller 2008) with Webb six-point
weights (B=4999, null imposed), clustered by source domain. ‡ Cluster-robust SEs
from rdrobust.\\
\end{minipage}
\end{table}

